\documentclass[11pt]{amsart}
\usepackage[T1]{fontenc}
\usepackage[margin=1.12in]{geometry}
\usepackage{amsmath,amssymb,amsthm,mathtools}
\usepackage{microtype}
\usepackage{xcolor}
\definecolor{linkblue}{RGB}{25,55,110}
\usepackage[colorlinks=true,linkcolor=linkblue,citecolor=linkblue,urlcolor=linkblue]{hyperref}
\usepackage{enumitem}
\usepackage{array}
\hypersetup{pdftitle={Uniqueness of rational solutions of framed tangle equations},
pdfauthor={Caleb Barnett}}
\setlist{itemsep=3pt,topsep=5pt}
\newtheorem{theorem}{Theorem}[section]
\newtheorem{proposition}[theorem]{Proposition}
\newtheorem{lemma}[theorem]{Lemma}
\newtheorem{corollary}[theorem]{Corollary}
\theoremstyle{definition}
\newtheorem{definition}[theorem]{Definition}
\newtheorem{convention}[theorem]{Convention}
\theoremstyle{remark}
\newtheorem{remark}[theorem]{Remark}
\newtheorem{example}[theorem]{Example}
\newcommand{\Z}{\mathbb Z}
\newcommand{\Q}{\mathbb Q}
\newcommand{\F}{\mathbb F}
\newcommand{\Sthree}{S^3}
\newcommand{\lCW}{\lambda_{\mathrm{CW}}}
\newcommand{\Tor}{\operatorname{Tor}}
\newcommand{\rk}{\operatorname{rank}}
\newcommand{\Det}{\operatorname{det}}
\newcommand{\Span}{\operatorname{span}}
\newcommand{\bd}{\partial}
\newcommand{\br}[1]{\langle #1\rangle}
\newcommand{\cl}[2]{C_{#1}(#2)}
\newcommand{\iso}{\cong_{\mathrm{or}}}
\newcommand{\tang}[1]{\langle #1\rangle}
\newcommand{\draftnote}[1]{\par\smallskip\noindent\textcolor{linkblue}{\textbf{Editorial placeholder.} #1}\par\smallskip}
\title[Rational solutions of tangle equations]{Uniqueness of rational solutions of framed tangle equations:\\ a proof of Sikora's conjecture and an unframed extension}
\author{Caleb Barnett}
\urladdr{https://barnettcaleb1.github.io}
\date{First complete draft for internal review: 3 November 2026}
\subjclass[2020]{57K10, 57K30, 57R58}
\keywords{tangle equation, rational tangle, double branched cover, Dehn filling, Casson--Walker invariant, Floer simple manifold}

\begin{document}
\begin{abstract}
We prove that a system of two balanced framed rational-closure equations,
with distinct underlying closing slopes, has a unique solution among
arbitrary two-string tangles whenever it has a rational solution.
This is Sikora's rational-solution uniqueness conjecture.
The proof uses uniqueness of the Kauffman bracket vector to identify the
integral peripheral homology of the double branched cover. Two actual
matched fillings then combine with the Casson--Walker surgery formula
and the Floer-simple manifold results of Rasmussen and Rasmussen to force
the cover to be a solid torus. We also prove an unframed version: two
actual closure identifications at distinct nonlongitudinal slopes,
together with equality of determinants at any third slope, suffice.
An elementary interpolation lemma recovers integral homology and the
longitude from these three scalar measurements. We give an arithmetic
condition under which two unframed probes suffice and describe the
remaining homological alternative explicitly.
\end{abstract}
\maketitle

\noindent\textbf{Draft status.}
This manuscript presents a deduction from published results for
adversarial review. It has not received external specialist review.
No claim of priority beyond the attribution discussion in
Section~\ref{sec:credit} is made.

\section{Introduction}\label{sec:intro}

A two-string tangle consists of two properly embedded arcs in a ball
with four fixed boundary points. A rational closing tangle \(P\)
assigns to an unknown tangle \(X\) a link
\(\cl P X=N(X+P)\). A system
\begin{equation}\label{eq:system}
 \cl P X=L_P,\qquad \cl Q X=L_Q
\end{equation}
asks how much two such closures determine about \(X\).
These equations arise in the topological study of DNA recombination;
the distinction between rational and arbitrary tangle solutions is
central to their use.

Sikora proves uniqueness among \emph{rational} solutions of a framed
system when \(P\) and \(Q\) have distinct unframed slopes
\cite[Proposition~1]{Sikora}. His conjecture asks for uniqueness among
all two-string tangles once a rational solution exists. The framing
convention is balanced framed isotopy, as in \cite[Section~4]{Sikora}.
Our first theorem proves that statement.

The mechanism has two different inputs. Equality of the two framed
closures determines the Kauffman bracket vector, by
\cite[Theorem~21]{Sikora}. This is a linear-algebraic statement, and by
itself does not identify the tangle. It gives determinant information
at auxiliary closure slopes. On the double branched cover, that
information determines integral homology and the preferred longitude.
The \emph{actual} closure identifications then supply matched
orientation-preserving Dehn fillings. These are the topological input
that polynomial identities alone would not supply.

For two finite-homology fillings, the Casson--Walker surgery formula
eliminates an unknown ambient invariant and forces the second
Alexander derivative to vanish. Two lens fillings also imply Floer
simplicity. The resulting alternating Alexander polynomial has
strictly positive second derivative unless it is constant.
The Thurston-norm formula of Rasmussen and Rasmussen then forces
boundary compression. A separate application of their generalized
solid-torus result treats a longitudinal probe.

The surgery--Alexander part of this argument has a close antecedent
in Gainullin's proof of \cite[Theorem~8.2]{Gainullin}. Here we use two
matched fillings and Floer simplicity of the exterior; we do not
assume that the auxiliary integral homology sphere is an \(L\)-space.
The relationship and the classical special cases are discussed in
Section~\ref{sec:credit}.

The unframed extension replaces the information recovered from the
bracket vector by three homology-order measurements. Orders of first
homology of torus-boundary fillings are absolute values of one
integral linear functional. Agreement with a primitive functional
at three distinct slopes forces agreement everywhere. Only two of
the three probes must be actual closure identifications. This also
gives an explicit arithmetic description of the ambiguity left by
two measurements.

Sikora proves that Jones unknot detection implies his framed
uniqueness conjecture \cite[Theorem~3]{Sikora}. We do not claim the
converse, and the argument here supplies no such implication.
In particular, a hypothetical nontrivial Jones-unit knot tied
locally into a rational tangle would preserve all bracket
measurements after compatible framing, while changing the actual
closure links. The surgery hypotheses would be missing.

The abstract of Sikora's IPAM talk of 15 October 2026 still states
the rational-solution assertion as a conjecture \cite{SikoraTalk}.
The literature checks underlying this draft located no prior proof
of the exact statement. This records the scope of those checks,
not a certification of priority or of the absence of unpublished work.

\section{Statements of the results}\label{sec:statements}

All balls and three-manifolds are oriented. Ambient link isotopies
preserve ambient orientation; link components need not be oriented.
Tangle comparisons retain the complete boundary marking, not only
the set of four endpoints. A slope on a torus is an unoriented
primitive class in its first integral homology. Distinct slopes are
therefore distinct elements of the projectivized integral lattice.

\begin{theorem}[Framed rational-solution uniqueness]\label{thm:framed}
Let \(P,Q\) be framed rational closing tangles with distinct
underlying unframed slopes, and let \(R\) be a framed rational
tangle. Suppose that \(X\) is a framed two-string tangle with the
same boundary marking and
\[
 \cl P X=\cl P R,\qquad \cl Q X=\cl Q R
\]
as balanced framed links. Then \(X=R\) as balanced framed tangles.
\end{theorem}

This is Conjecture~2 of the specified arXiv version of
\cite{Sikora}, numbered Conjecture~1.2 in the journal article.
An arbitrary collection of closed components may also be allowed
in \(X\): the proof excludes them from the hypotheses.
The two closing tangles must differ after forgetting framing;
two framings of one closing slope do not suffice.

For a compact oriented manifold \(M\) with one torus boundary, write
\(M(\alpha)\) for Dehn filling along \(\alpha\), and put
\[
 h(Z)=
 \begin{cases}|H_1(Z;\Z)|,&H_1(Z;\Z)\text{ finite},\\
 0,&H_1(Z;\Z)\text{ infinite}.
 \end{cases}
\]
A \emph{marked solid torus} \(V\) is a solid torus whose oriented
boundary slope lattice has been identified with that of \(M\).
Its disk slope will be denoted \(\lambda\).

\begin{theorem}[Homology interpolation]\label{thm:homology}
Let \(M\) be compact, connected and oriented with one torus boundary,
and let \(V\) be a marked solid torus. Suppose that \(M\) has a
rational-homology-sphere filling. If
\[
 h(M(\alpha_i))=h(V(\alpha_i)),\qquad i=1,2,3,
\]
at three distinct slopes, then \(H_1(M;\Z)\cong\Z\), the
boundary-to-homology map is surjective, and its kernel is generated
by the disk slope of \(V\).
\end{theorem}

For a rational tangle \(R\), a closing slope is
\emph{nonlongitudinal} when it is different from the disk slope of
its double branched cover. Equivalently, the corresponding closure
of \(R\) has nonzero determinant.

\begin{theorem}[Unframed three-probe rigidity]\label{thm:three}
Let \(R\) be rational and \(X\) be any two-string tangle in the same
marked ball. Let \(\alpha,\beta\) be distinct nonlongitudinal
rational closing slopes for \(R\). Suppose
\[
 \cl\alpha X\simeq\cl\alpha R,\qquad
 \cl\beta X\simeq\cl\beta R
\]
as unframed links by ambient isotopies. If, at one further slope
\(\gamma\notin\{\alpha,\beta\}\),
\[
 \Det\cl\gamma X=\Det\cl\gamma R,
\]
then \(X\simeq R\) as unframed marked tangles.
\end{theorem}

\begin{corollary}[Three actual closures]\label{cor:actual}
If three distinct rational closures of \(X\) are ambient isotopic
to the corresponding closures of a rational tangle \(R\), then
\(X\simeq R\) as unframed marked tangles.
\end{corollary}

\begin{corollary}[Diagnostic form]\label{cor:diagnostic}
Suppose \(X\not\simeq R\), with \(R\) rational, and two actual
unframed closure equalities hold at distinct nonlongitudinal slopes
\(\alpha,\beta\). Then the determinants of the corresponding
closures differ at every other slope.
\end{corollary}

Let \(\Delta(\alpha,\beta)\) denote slope distance, the absolute
value of the determinant of their primitive representatives.

\begin{theorem}[Arithmetic two-probe criterion]\label{thm:arithmetic}
Under the two actual closure hypotheses of
Theorem~\ref{thm:three}, put \(d=\Delta(\alpha,\beta)>0\).
If
\[
 d\nmid 2\Det\cl\alpha R,
\]
then \(X\simeq R\), without any third-probe assumption.
The divisibility condition is unchanged if \(\beta\) replaces
\(\alpha\).
\end{theorem}

The auxiliary surgery statement used in both tangle theorems is
worth isolating.

\begin{proposition}[Two marked finite fillings]\label{prop:surgery}
Let \(M\) be compact, connected and oriented with one torus
boundary, and let \(V\) be a marked solid torus. Assume that
\(H_1(M;\Z)=\Z\), the boundary map is surjective, and the
homological longitude of \(M\) agrees with the disk slope of \(V\).
If two distinct nonlongitudinal slopes satisfy
\[
 M(\alpha)\iso V(\alpha),\qquad M(\beta)\iso V(\beta),
\]
then \(M\) is a solid torus with the same disk slope as \(V\).
\end{proposition}

The marking in this proposition is indispensable. The
homeomorphisms of the filled manifolds need not respect their
filling cores. The same signed surgery coefficients are determined
by the marking before those homeomorphisms are used.

\section{Conventions and the inputs from the literature}\label{sec:prelim}

\subsection{Balanced framing and the bracket}

Fix four named endpoints on the boundary of an oriented ball.
We use blackboard-framed diagrams and Sikora's balanced
Reidemeister equivalence \cite[Section~4 and Figure~5]{Sikora}.
It includes the usual second and third moves and the balanced
creation or removal of opposite curls. In particular, compensating
framing changes may be moved between components. This is more
flexible than isotopy of independently framed components.
An isolated positive framing twist multiplies the bracket by
\(-A^3\). The total framing information relevant to the bracket
is not discarded.

Write \(\mathcal R=\Z[A,A^{-1}]\) and
\(\delta=-A^2-A^{-2}\). The two trivial boundary pairings
\(\tang0,\tang\infty\) form the basis of the relative bracket
module of the marked ball:
\[
 [X]=x_0[\tang0]+x_\infty[\tang\infty],
 \qquad x_0,x_\infty\in\mathcal R.
\]
This vector is invariant under balanced framed isotopy.
Closed components in \(X\), if allowed, are evaluated by the same
skein relations and do not change the rank of this module.

For nonempty closed links we use the reduced bracket
\(\br U=1\), so adding a disjoint trivial component multiplies
the bracket by \(\delta\). If one starts instead with the
absolute convention \(\br{\varnothing}_{\rm abs}=1\),
then \(\br L=\br L_{\rm abs}/\delta\).
This division is a polynomial division made \emph{before}
any evaluation at a zero of \(\delta\).

For a fixed closing tangle \(T\), closure is linear:
\begin{equation}\label{eq:closurefunctional}
 \br{\cl T X}
 =x_0\br{N(\tang0+T)}+
  x_\infty\br{N(\tang\infty+T)}.
\end{equation}
If \([T]=(a,b)^{\mathsf T}\), the row vector of this
functional is
\[
 (\delta a+b,\ a+\delta b).
\]
For two probes its determinant is
\begin{equation}\label{eq:closurematrix}
 (\delta^2-1)(a_1b_2-a_2b_1).
\end{equation}
For distinct rational probe slopes the last factor is nonzero.
This is the independence assertion in the proof of Sikora's
Theorem~21; the rational bracket ratio recovers the rational
fraction at its classical specialization
\cite[Corollary~19 and Theorem~21]{Sikora}.
Consequently, two equal framed closure brackets force equality
of the two relative bracket vectors, over \(\Q(A)\) and hence
over \(\mathcal R\).

We will also use \cite[Corollary~20]{Sikora} in precisely its
rational-only form: equal bracket vectors of two balanced framed
rational tangles imply equality of the tangles.
It makes no assertion here about detecting a rational tangle
among all tangles.

\subsection{Branched covers and their markings}

The connected double cover of the boundary sphere, branched at
its four marked points, is a torus. The cover of a
rational tangle is a solid torus. Attaching a rational closing
tangle is Dehn filling along its disk slope upstairs. Thus, if
\(M=\Sigma_2(B,X)\), then
\begin{equation}\label{eq:coverclosure}
 \Sigma_2(\cl\alpha X)=M(\alpha).
\end{equation}
The covering orientation is the one for which the local covering
map preserves orientation off the branch set.

The fixed marking identifies the two boundary branched covers
over the \emph{same} marked sphere. Its two possible lifts differ
by the deck involution. In the pillowcase model
\(T^2/(z\mapsto-z)\), this involution acts as \(-I\) on
\(H_1(T^2;\Z)\). Therefore the lifted slope identification is
well defined on unoriented slopes. It also preserves the boundary
orientation.

For later signed coefficients, give the peripheral torus the
orientation as the boundary of a tubular neighborhood of the
filling core. This is the negative of its orientation as the
boundary of the exterior. We choose oriented representatives
\(\mu,\lambda\) with
\(\mu\cdot\lambda=+1\) in this peripheral convention.
Using the exterior orientation instead reverses both
intersection conventions at once; it does not allow an
independent sign change on one side of a comparison.

An ambient orientation-preserving isotopy of branch links lifts
to an orientation-preserving homeomorphism of the closed covers.
Away from the branch set this follows from the covering
orientation, and extension over a branched tubular neighborhood
preserves it. The lifted homeomorphism can forget both the
marked boundary and the filling core. Neither must be preserved
to equate invariants of the closed oriented manifolds.

We use the classical rationality criterion in the form stated by
Buck and Mauricio \cite[Theorem~3.5(2)]{BM}, attributed there to
Lickorish \cite{Lickorish}: a two-string tangle has solid-torus
double branched cover if and only if it is rational.
This statement applies to the given branched cover and its
involution. It is not a classification obtained by changing the
involution. Rational tangles with the fixed four-point marking
are classified by their slopes
\cite[Theorem~3.6]{BM}. Equivalently, their disk slopes in the
marked boundary torus classify them.

\subsection{Determinants and homology}

Put \(\zeta=e^{\pi i/4}\). The determinant of a link is
\[
 \Det L=|\br L(\zeta)|
       =h(\Sigma_2(L)).
\]
The framing factor has absolute value one at \(\zeta\).
These equalities hold also when the determinant is zero, with
the convention defining \(h\).
Here is a presentation argument that includes the zero case.
Arrange a connected projection of the link; if necessary,
Reidemeister-II moves between projected components accomplish this
without changing the link. Let \(G\) be a reduced Goeritz matrix.
Lickorish's Theorem~9.3 identifies the integral form of a connected
spanning surface as a presentation of \(H_1(\Sigma_2(L);\Z)\);
Theorem~9.4 identifies the checkerboard-surface form with \(G\),
and Corollary~9.5 gives \(\Det L=|\det G|\)
\cite[Chapter~9, pp.~96--99]{LickorishBook}.
Thus the finite cokernel has order \(|\det G|\) when \(G\) is
nonsingular, and its cokernel has positive rank when \(G\) is
singular. The \(0\times0\) matrix presents the trivial group
and has determinant one. This proves both cases of the displayed
identity. See also \cite[Section~3]{DFKLS} for the
determinant--bracket formulation.

We need only the following consequence of
\cite[Proposition~6.5]{Stoimenow}: a link has odd determinant
if and only if it is a knot. In particular, a determinant-one
closure of a two-string tangle is a knot. If the tangle had
a closed component, that component would survive into every
closure alongside at least one component containing its
proper arcs. Thus a determinant-one closure excludes closed
components of the tangle.

\subsection{The signed surgery formula}

We normalize the Alexander polynomial of a knot \(K\) in an
integral homology sphere \(Y\) by
\[
 \Delta_K(t^{-1})=\Delta_K(t),\qquad \Delta_K(1)=1.
\]
Write \(Y_{p/q}(K)\) for filling its exterior along
\(p\mu+q\lambda_K\), with \(\lambda_K\) the preferred
longitude. Our lens-space convention is
\(L(p,q)=S^3_{p/q}(U)\), including its orientation.

Ni and Wu \cite[Proposition~3.1, equation~(7)]{NiWu}
state the Casson--Walker surgery formula for a
null-homologous knot in a rational homology sphere.
In the integral-homology-sphere case used here it is
\begin{equation}\label{eq:cw}
 \lCW(Y_{p/q}(K))
 =\lCW(Y)+\lCW(L(p,q))
       +\frac{q}{2p}\Delta_K''(1),\qquad p\ne0.
\end{equation}
We cite that explicitly checked formulation, after
Boyer and Lines \cite{BoyerLines}. No assumption that \(Y\)
is \(S^3\), or that \(Y\) is an \(L\)-space, is present.
Changing both \(p,q\) to their negatives changes neither
the slope nor the coefficient \(q/(2p)\).
Reversing the orientation of just one filled manifold
would change its Casson--Walker invariant, so it is not
permitted in the comparisons below. The surgery fraction \(p/q\)
is the coordinate of the lifted filling slope, and need not equal
the Conway fraction naming the closing tangle downstairs. The same
fixed closing tangle supplies the same lifted slope on both covers.

\subsection{The Floer-simple input}\label{subsec:floer}

We state exactly the portions of Rasmussen and Rasmussen
\cite{RR} used here; numbering refers to arXiv:1508.05900v2.
An \(L\)-space is a rational homology sphere with minimal-rank
Heegaard Floer homology. Lens spaces, including \(S^3\), are
\(L\)-spaces in either orientation.

For a connected oriented manifold \(M\) with one torus
boundary and \(b_1(M)=1\), their Proposition~1.3 says that
having more than one \(L\)-space filling is equivalent to
Floer simplicity. Their torsion normalization preceding
Proposition~1.4, together with Section~2.1, Proposition~1.4
and Proposition~2.2, has the following consequence when \(M\) is Floer simple and
\(H_1(M;\Z)=\Z\):
\begin{equation}\label{eq:binarytorsion}
 \tau(M)=\frac{D(t)}{1-t}
         =\sum_{n\ge0}a_nt^n,\qquad
 a_0=1,\quad a_n\in\{0,1\},\quad a_n=1\ (n\gg0).
\end{equation}
The choice of generator and a power of \(t\) fixes the
normalization. If \(M\) is a knot exterior in an integral
homology sphere, then \(D(t)=t^h\Delta_K(t)\), with the
positive sign fixed by \(D(1)=\Delta_K(1)=1\).
Here \(h\) is a shift, not an assumed genus.

More generally, \(\overline\Delta(M)\) is the polynomial
obtained by projecting the Turaev torsion to the free
part of homology and multiplying by \(1-t\).
When \(H_1(M)=\Z\) the projection is the identity, so
\begin{equation}\label{eq:barD}
 \overline\Delta(M)=D(t).
\end{equation}
We use degree to mean Laurent breadth:
\(\deg f=\max\operatorname{supp}(f)-\min\operatorname{supp}(f)\).
For \(D\) in \eqref{eq:binarytorsion}, its constant term
is one and its lowest exponent is zero, so this is its
ordinary polynomial degree.

The first assertion of \cite[Corollary~2.3]{RR} says that
if \(M\) is boundary incompressible and Floer simple, then
\begin{equation}\label{eq:norm}
 \|M\|=\deg\overline\Delta(M)-1,
\end{equation}
where \(\|M\|\) is the Thurston norm of a generator of
\(H_2(M,\bd M;\Z)\). We will have irreducibility as well,
although the stated norm assertion does not require it.
We do not use the separate fibering assertion of that corollary.

Finally \cite[Proposition~7.8]{RR} states that, for a rational
homology sphere \(Z\), a knot in \(Z\#(S^1\times S^2)\)
admitting an \(L\)-space surgery has
generalized-solid-torus exterior. We use only \(Z=S^3\).
By Definition~7.2, such an exterior is Floer simple and
\[
 \deg\overline\Delta(M)<g_M.
\]
Here \(g_M\) is the order of the image of the homological
longitude in the torsion subgroup of \(H_1(M;\Z)\).
The order of its zero element is one. In the torsion-free
case \(g_M=1\); equivalently the relation
\(|\Tor H_1(M)|=k_Mg_M\) in \cite[Section~7, introductory notation]{RR}
forces both positive integers \(k_M,g_M\) to be one.
Generalized solid tori need not be actual solid tori.
We will use the torsion-free hypothesis and irreducibility
before making that conclusion.

\section{Proof of the framed theorem}\label{sec:framedproof}

We give the proof in five steps and isolate the common
surgery argument for use in the unframed theorem.

\subsection*{Step 1. Equality of relative brackets}

The hypotheses of Theorem~\ref{thm:framed} give two
\emph{exact} reduced bracket equalities: the closing
tangle and its framing are the same on both sides of
each equation. Apply the two independent functionals
\eqref{eq:closurefunctional} to \([X]-[R]\).
Their determinant \eqref{eq:closurematrix} is nonzero
because the probe slopes differ. Thus
\begin{equation}\label{eq:vector}
 [X]=[R]\quad\text{in }\mathcal R^2.
\end{equation}
This is precisely the conclusion of
\cite[Theorem~21]{Sikora}; the displayed linear proof
also applies if \(X\) initially has closed components.
Every additional rational closing tangle therefore
gives equal brackets, and hence equal determinants.
It does not yet give isotopic links at those
additional slopes.

\subsection*{Step 2. Integral peripheral data and irreducibility}

Let \(M=\Sigma_2(B,X)\) and \(V=\Sigma_2(B,R)\).
The latter is a solid torus. Let \(\lambda\) be its
primitive disk slope. Choose a slope \(\mu\) with
\(\mu\cdot\lambda=1\), different from both original
probes. Such a choice can be made before any filling:
the slopes \(\mu_0+n\lambda\), \(n\in\Z\), give
infinitely many choices. Orient them compatibly with
the peripheral convention.

The corresponding rational closures of \(R\) are
the two-component unlink and an unknot. Upstairs,
\[
 V(\lambda)=S^1\times S^2,\qquad V(\mu)=S^3.
\]
Equation~\eqref{eq:vector} therefore gives
\begin{equation}\label{eq:01}
 h(M(\lambda))=0,\qquad h(M(\mu))=1.
\end{equation}
The determinant-one closure of \(X\) is a knot,
so \(X\) has no closed components.

Set \(Y=M(\mu)\). By \eqref{eq:01}, this is an integral
homology sphere, and \(M\) is the exterior of the
core \(K\) of the attached solid torus \(W\).
For completeness, Mayer--Vietoris gives
\[
 0=H_2(Y)\longrightarrow H_1(\bd M)
 \longrightarrow H_1(M)\oplus H_1(W)
 \longrightarrow H_1(Y)=0.
\]
Since \(H_1(\bd M)=\Z^2\) and \(H_1(W)=\Z\), it follows
that \(H_1(M)=\Z\). The preceding part of the sequence,
or Euler characteristic and duality, gives \(H_2(M)=0\).
The meridian \(\mu\) maps to a generator of \(H_1(M)\).
The kernel of the boundary map is generated by the
primitive preferred longitude \(\lambda_K\).

Filling along \(p\mu+q\lambda_K\) has first homology
\(\Z/p\Z\), where \(p=0\) means \(\Z\).
The first equality of \eqref{eq:01} therefore forces
\(\lambda=\pm\lambda_K\). These are the same
unoriented slope. We have identified a common
\emph{integral} meridian--longitude pair for \(M\)
and \(V\), not merely a rational longitude.
No identification of \(Y\) with \(S^3\) has been made.

At least one original probe is different from
\(\lambda\). At that slope, the actual closure
isotopy gives an actual finite-homology lens-space
filling of \(M\).

\begin{lemma}\label{lem:irreducible}
If \(M\) has one torus boundary, \(H_1(M;\Z)=\Z\),
and a finite-homology lens-space filling, then \(M\)
is irreducible.
\end{lemma}
\begin{proof}
In a prime decomposition write \(M=M_0\# Z_1\#\cdots\#Z_s\),
where \(M_0\) carries the boundary and the \(Z_j\) are
closed prime factors. For any compact oriented
manifold with one torus boundary,
\(\chi(M_0)=0\), so \(b_1(M_0)=1+b_2(M_0)\ge1\).
The identity \(H_1(M)=\Z\) forces
\(H_1(M_0)=\Z\) and \(H_1(Z_j)=0\).
In particular no \(S^1\times S^2\) factor is possible.
After filling,
\[
 M(\alpha)=M_0(\alpha)\#Z_1\#\cdots\#Z_s.
 \]
A finite-homology lens space is prime. A closed
homology-sphere summand \(Z_j\) must be \(S^3\):
if it were the nontrivial prime factor of the
filling, that filling itself would be a
homology-sphere lens space, hence \(S^3\).
Prime decomposition therefore removes all closed
factors. A prime orientable manifold with nonempty
boundary is irreducible; the closed exception
\(S^1\times S^2\) does not arise.
\end{proof}

\subsection*{Step 3. Two nonlongitudinal probes}

Suppose first that neither original probe is
\(\lambda\). Write
\(\alpha_i=p_i\mu+q_i\lambda\), where
\(\gcd(p_i,q_i)=1\) and \(p_i\ne0\). The choice of
\(\mu\) away from the probes also ensures \(q_i\ne0\).
We spell out the sign issue before applying
\eqref{eq:cw}.

\begin{lemma}[Signed comparison]\label{lem:signs}
With the common marking and the meridian and
longitude fixed as above, the actual closure
isotopies give
\begin{equation}\label{eq:signed}
 Y_{p_i/q_i}(K)\iso S^3_{p_i/q_i}(U),\qquad i=1,2,
\end{equation}
with identical signed coefficients on both sides.
\end{lemma}
\begin{proof}
The boundary identification is lifted from the
identity of the marked four-punctured sphere.
Its two choices differ by the deck action \(-I\).
They replace both \(\mu,\lambda\) by
\(-\mu,-\lambda\), and hence \((p_i,q_i)\) by
\((-p_i,-q_i)\). The quotient \(p_i/q_i\)
does not change.

The chosen slope \(\mu\) is the meridian of
both filling cores, since both are defined by
\(\mu\)-filling. The determinant-zero argument
has already identified their preferred
longitudes with the same slope \(\lambda\).
Orient the pairs so that their intersection is
\(+1\) in the common peripheral convention. Orient each core
so that \(\lambda\) runs once positively along it; then \(\mu\)
is its positive meridian.
Reversing an oriented knot reverses both
peripheral classes; it does not change the
coefficient. There is no independent replacement
of a longitude by its negative, or by a
longitude plus a meridian: the former would
reverse the peripheral orientation without
also reversing the meridian, and the latter
would no longer be null-homologous.
Likewise a meridional shear would change
the chosen filling slope \(\mu\).

The core in \(V(\mu)=S^3\) is an unknot.
Thus the marked slope \(\alpha_i\) represents
the same signed surgery on each core.
Finally the original link isotopy lifts to
an orientation-preserving homeomorphism of
the closed branched covers. This establishes
\eqref{eq:signed}. The lifted homeomorphism
need not carry one core to the other: it is
used only after the coefficients have been
identified independently.
\end{proof}

Applying \eqref{eq:cw} to \eqref{eq:signed} and
cancelling the common lens-space term yields
\begin{equation}\label{eq:twoCW}
 0=\lCW(Y)+\frac{q_i}{2p_i}\Delta_K''(1),
 \qquad i=1,2.
\end{equation}
Distinct slopes with \(p_i\ne0\) have distinct
ratios \(q_i/p_i\). Subtracting gives
\begin{equation}\label{eq:zeroderivative}
 \Delta_K''(1)=0,\qquad \lCW(Y)=0.
\end{equation}
The second equality does not identify \(Y\).
The use of surgery comparison and the
Alexander derivative follows the strategy
of \cite[Theorem~8.2 and its proof]{Gainullin};
here the second equation removes the
unknown \(\lCW(Y)\).

Both fillings are \(L\)-spaces. Their distinct filling slopes
and \cite[Proposition~1.3]{RR} make \(M\)
Floer simple. Since \(H_1(M)=\Z\),
\eqref{eq:binarytorsion} applies. The
coefficients of \(D=(1-t)\tau(M)\) are
\(a_n-a_{n-1}\), with \(a_{-1}=0\).
The nonzero differences of a binary sequence
alternate in sign, start with \(+1\), and
end with \(+1\), because the sequence is
eventually one. Also \(D(1)=1\).

Symmetry of \(\Delta_K\) now implies that
for some \(k\ge0\),
\begin{equation}\label{eq:alternating}
 \Delta_K(t)=(-1)^k+
   \sum_{j=1}^k(-1)^{k-j}(t^{n_j}+t^{-n_j}),
 \qquad 0<n_1<\cdots<n_k .
\end{equation}
Indeed the number of nonzero coefficients
is odd, and all noncentral coefficients
occur in symmetric pairs. A central
coefficient must therefore be present,
with the sign shown.
Differentiating twice gives
\begin{equation}\label{eq:positive}
 \Delta_K''(1)=
 2\sum_{j=1}^k(-1)^{k-j}n_j^2.
\end{equation}
For even \(k>0\), the half-sum is
\((n_k^2-n_{k-1}^2)+\cdots+(n_2^2-n_1^2)>0\).
For odd \(k>0\), it is the same sum of
positive successive differences with
the additional term \(n_1^2>0\).
Thus \eqref{eq:zeroderivative} forces \(k=0\)
and \(\Delta_K=1\).
Because \(D\) has constant term one,
its shift is now zero: \(D=1\).

If \(\bd M\) were incompressible,
\eqref{eq:barD} and \eqref{eq:norm} would give
\(\|M\|=0-1=-1\), contrary to nonnegativity
of the Thurston norm. Hence the boundary
compresses.

\begin{lemma}\label{lem:compress}
An irreducible compact orientable
three-manifold with a single torus boundary
and a boundary-compressing disk is a solid torus.
\end{lemma}
\begin{proof}
The boundary of the compressing disk is
essential and nonseparating on the boundary
torus. A regular neighborhood of the
torus together with the disk has, besides
the original torus, a sphere boundary.
Irreducibility fills that sphere on the
interior side by a ball. Reversing the
compression attaches one one-handle to
that ball, giving a solid torus. The sphere
cannot bound its ball on the side
containing the original boundary torus,
so there is no additional component left.
\end{proof}

Lemmas~\ref{lem:irreducible} and
\ref{lem:compress} show that \(M\) is a
solid torus. This also proves
Proposition~\ref{prop:surgery}: its
homology hypotheses give the same choice
of \(Y,K,\mu,\lambda\), and its two
oriented filling identifications are
exactly \eqref{eq:signed}. A solid torus's
disk slope is its homological longitude,
so the final slope agrees with that of \(V\).

\subsection*{Step 4. A longitudinal probe}

If one original probe is \(\lambda\),
the actual closure equality gives
\[
 M(\lambda)=S^1\times S^2,
\]
not merely a first-homology isomorphism.
The other probe is nonlongitudinal and
gives an \(L\)-space lens filling.
Regard \(M\) as the exterior of the
\(\lambda\)-filling core in \(S^1\times S^2\).
Apply \cite[Proposition~7.8]{RR} with
\(Z=S^3\) and with this other filling as
the stipulated \(L\)-space surgery.
It follows that \(M\) is a generalized
solid torus.

Step~2 has already proved \(H_1(M)=\Z\).
Therefore \(g_M=1\), both directly from
the zero image of the longitude and
from \(|\Tor H_1(M)|=k_Mg_M=1\).
Definition~7.2 of \cite{RR} gives
Floer simplicity and
\(\deg\overline\Delta(M)<1\).
This normalized polynomial is nonzero;
in the torsion-free case its value at
one is one. Its degree is therefore zero.
Boundary incompressibility would again
give \(\|M\|=-1\) by Corollary~2.3.
The boundary compresses, and the
irreducibility already proved in Step~2
allows Lemma~\ref{lem:compress} to
conclude that \(M\) is a solid torus.

This step uses Proposition~7.8 at arbitrary
distance from the longitudinal slope.
Neither Property \(R\) in \(S^3\) alone
nor an assertion that every generalized
solid torus is a solid torus would suffice.
No uniqueness assertion for a single
unlink closure is used.

\subsection*{Step 5. The marked balanced framed tangle}

In both cases the double branched cover
of the given two-arc tangle \(X\) is a
solid torus. By \cite[Theorem~3.5(2)]{BM},
\(X\) is rational. This criterion is
about the actual branched covering
involution. Keeping the fixed
boundary parametrization then places
\(X\) in the usual marked rational-tangle
classification. Extending a boundary
parametrization over a ball uses the
Alexander trick; no change of boundary
marking is being inserted in the
comparison with \(R\).

We now have two balanced framed
\emph{rational} tangles with the same
bracket vector \eqref{eq:vector}.
Sikora's Corollary~20 applies and gives
\(X=R\) as balanced framed tangles.
More explicitly, the bracket ratio
first determines the common rational
fraction. Any remaining balanced
framing discrepancy is an integer \(n\)
and multiplies the nonzero bracket
vector by \((-A^3)^n\); equality forces
\(n=0\). This completes the proof of
Theorem~\ref{thm:framed}.

\section{The unframed extension}\label{sec:unframed}

\subsection{Homology orders as a linear functional}

We prove Theorem~\ref{thm:homology} without
any skein or Floer-theoretic input.
A finite-homology filling implies
\(b_1(M)\le1\), since filling kills at
most one generator of rational first
homology. On the other hand
\(\chi(M)=\frac12\chi(\bd M)=0\) gives
\(b_1(M)=1+b_2(M)\ge1\).
Hence \(b_1(M)=1\). Write
\[
 H_1(M;\Z)=\Z\oplus T,\qquad t_0=|T|.
\]
Let \(\ell_M:H_1(\bd M;\Z)\to\Z\)
be the free projection of the
inclusion-induced homomorphism.
It is not assumed primitive.

For a primitive slope \(\alpha\),
Mayer--Vietoris for filling identifies
\[
 H_1(M(\alpha);\Z)=
 H_1(M;\Z)/\langle i_*(\alpha)\rangle .
\]
To see this quotient directly, a
complementary boundary curve maps to
the generator of the attached solid
torus; eliminating that generator
leaves exactly the relation
\(i_*(\alpha)=0\).
Write \(i_*(\alpha)=(n,t)\).
If \(n=0\), the quotient still has
rank one. If \(n\ne0\), its order is
\(|n|\,|T|\): reduction of the free
coordinate modulo \(n\) gives a
surjection to \(\Z/n\Z\), whose kernel
has order \(|T|\), since the cyclic
relation has no nonzero element in
\(0\oplus T\).
Thus for every slope
\begin{equation}\label{eq:linearh}
 h(M(\alpha))=t_0\,|\ell_M(\alpha)|.
\end{equation}

Let \(l:H_1(\bd V;\Z)\to\Z\) be the
primitive homology functional of the
solid torus and use the marking to
regard it as a functional on
\(H_1(\bd M)\). Set \(m=t_0\ell_M\).
The three hypotheses say
\(m(\alpha_i)^2=l(\alpha_i)^2\).
The homogeneous quadratic polynomial
\[
 m^2-l^2=(m-l)(m+l)
\]
therefore vanishes on three distinct
rational directions. A nonzero linear
form vanishes on only one such
direction. By the pigeonhole principle
one of \(m-l,m+l\) vanishes on two
independent vectors, and is zero.
Consequently \(m=l\) or \(m=-l\).
Because \(l\) is primitive integral
and \(m=t_0\ell_M\), the positive
integer \(t_0\) is one and \(\ell_M\)
is primitive. Its kernel is the disk
slope of \(V\). This proves
Theorem~\ref{thm:homology}.

In fact the explicit assumption of
a finite filling is redundant once
the three equalities hold: a solid
torus has only one zero-order filling
slope, so at least two of the three
orders are finite and nonzero.
We retained the assumption to state
the general linear-order formula
in its natural setting.

\subsection{Two isotopies and one determinant}

For Theorem~\ref{thm:three}, take the
double branched covers \(M,V\) of
\(X,R\). The actual isotopies at
\(\alpha,\beta\) give
\[
 M(\alpha)\iso V(\alpha),\qquad
 M(\beta)\iso V(\beta).
\]
These are finite lens fillings by
the nonlongitudinal assumption.
Together with the determinant
equality at \(\gamma\), they supply
the three homology-order equalities.
Theorem~\ref{thm:homology} gives
\(H_1(M)=\Z\), surjectivity of the
boundary map, and the common
longitude. Proposition~\ref{prop:surgery}
now makes \(M\) a solid torus with
the marked disk slope of \(V\).

If closed components were initially
allowed in \(X\), the recovered
primitive boundary map provides a
slope \(\mu\) with \(h(M(\mu))=1\).
Its closure is a knot, so again
there are no closed components.
The rationality criterion applies.
Rational-tangle classification by
the marked disk slope gives
\(X\simeq R\), now without a framing
assertion. This proves the theorem.
The route uses neither the
Kauffman bracket-vector theorem nor
the generalized-solid-torus result.

For Corollary~\ref{cor:actual}, at
most one of three slopes is
longitudinal. Use the other two as
the actual finite probes and use
the remaining equality only for
its determinant. For
Corollary~\ref{cor:diagnostic}, any
third determinant equality would
contradict \(X\not\simeq R\) by
Theorem~\ref{thm:three}. The
diagnostic applies to every
nonidentical solution, including a
different rational one.

\subsection{The two-probe arithmetic and its residual alternative}

We prove Theorem~\ref{thm:arithmetic}
and give the exact alternative
when its divisibility test fails.
There is a finite filling, so
the notation \(m=t_0\ell_M\)
and \(l\) from \eqref{eq:linearh}
is available. Choose primitive
representatives and an oriented
integral basis in which
\[
 \alpha=(1,0),\qquad
 \beta=(a,d),\qquad d>0,\quad
 \gcd(a,d)=1.
\]
Write \(l=(u,v)\), so \(u\ne0\)
and \(|u|=\Det\cl\alpha R\).
Changing the global sign of \(m\)
if necessary, the first order
equality gives \(m(\alpha)=u\).
Write \(m=(u,w)\). The second
order equality is
\[
 |au+dw|=|au+dv|.
\]
There are exactly two possibilities:
\begin{equation}\label{eq:reflected}
 m=l,\qquad\text{or}\qquad
 m=m^*:=\left(u,-v-\frac{2au}{d}\right).
\end{equation}
In the first case, primitivity of
\(l=t_0\ell_M\) forces torsion-free
homology and a primitive boundary
map, with the correct longitude.
Proposition~\ref{prop:surgery}
finishes.

The second case is integral only
if \(d\mid2au\), equivalently
\(d\mid2u\), since \(\gcd(a,d)=1\).
The hypothesis excludes this case.
Finally \(l(\beta)=au+dv\) and
\(\gcd(a,d)=1\) show
\[
 d\mid2l(\beta)\quad\Longleftrightarrow\quad
 d\mid2u.
\]
This proves symmetry and the
arithmetic theorem.

Equation~\eqref{eq:reflected} gives
more than a divisibility obstruction.
If a nonidentical solution exists,
then up to global sign its complete
homology-order functional is
\(m^*\). Thus
\[
 \Det\cl\gamma X=|m^*(\gamma)|
\]
for every closing slope \(\gamma\).
The kernel of \(m^*\) determines
the rational longitude, and the
greatest common divisor of its
two coordinates is
\(|\Tor H_1(M)|\) times the index
of the image of \(\ell_M\).
This product is determined by the
two-probe data; its two factors
are not thereby separated.

If one closure is an unknot,
\(|u|=1\), so distance at least
three satisfies the criterion.
Distance two always passes the
divisibility test and is not
settled by this arithmetic
argument. No new surgery
characterization is claimed for
the unknot-closure special case.

\subsection{An exact two-probe calibration}\label{subsec:ES}

The following calibration is the mirror of the first two solutions
in \cite[Example~3, pp.~33--34]{ES1999}, followed by integer
transfers in the Montesinos presentation. In our rational-tangle
convention, take
\[
 X=\tang{-1}+\tang{1/2}+\tang{1/3},
 \qquad R=\tang{1/4}.
\]
The numerator closures at add-ons
zero and one agree: they are
respectively an unknot and the
same cinquefoil, with handedness
fixed by the rational-tangle
convention. The tangle \(X\)
is not rational: its cover is
Seifert fibred over a disk with
exceptional orders two and three,
the trefoil exterior up to
orientation, and is not a solid
torus.

For an add-on \(\tang{p/q}\),
the Montesinos determinant formula
gives
\begin{align}
 \Det N(X+\tang{p/q})
 &=\left|6q\left(-1+\frac12+\frac13+\frac pq\right)\right|
   =|6p-q|,\label{eq:ESX}\\
 \Det N(R+\tang{p/q})&=|4p+q|.\label{eq:ESR}
\end{align}
These formulas include \(q=0\)
by the boundary homology
functionals, or by the
denominator closure. Their
squared difference is
\begin{equation}\label{eq:ESdiff}
 (6p-q)^2-(4p+q)^2=20p(p-q).
\end{equation}
For primitive \((p,q)\) the
determinants therefore agree
only at the add-ons \(0\) and
\(1\). Every third determinant
separates these two tangles.
For example, the denominator
closure determinants are six
and four. The two common
probe slopes have distance
one, exactly a case not
excluded by
Theorem~\ref{thm:arithmetic}.
This example shows why the
third measurement in
Theorem~\ref{thm:three}
cannot be omitted uniformly.

\section{Stress tests and the role of the hypotheses}\label{sec:stress}

The following examples identify
which steps would fail under
natural weakenings. They are
not counterexamples to the
theorems as stated.

\subsection{An \texorpdfstring{\(L\)}{L}-space filling is not an actual lens filling}

Let \(Z\) be a nontrivial
integral homology sphere that
is an \(L\)-space, for example
the Poincare homology sphere,
and set
\[
 M=(S^1\times D^2)\# Z.
\]
Then \(H_1(M)=\Z\), the
longitude is primitive,
and \(M\) has many \(L\)-space
fillings \(L(p,q)\# Z\).
Its torsion normalization is
the one with
\(\overline\Delta(M)=1\);
compare \cite[Section~2.1]{RR}.
Nevertheless \(M\) is not a
solid torus. The norm
argument allows boundary
compression, but it does not
remove the closed summand.
The actual lens-space filling
in Lemma~\ref{lem:irreducible}
is what removes it.

\subsection{Floer simplicity does not force a constant Alexander polynomial}

The trefoil exterior is
irreducible, has
\(H_1=\Z\), and admits
multiple \(L\)-space
fillings. Its symmetric
Alexander polynomial is
\[
 \Delta(t)=t-1+t^{-1},
 \qquad \Delta''(1)=2.
\]
It therefore has precisely
the binary torsion and
alternating coefficients
used in Step~3, without
being a solid torus.
The missing condition is
the pair of correctly
marked unknot-surgery
comparisons. Their
Casson--Walker subtraction,
not Floer simplicity alone,
forces the derivative to
vanish.

More generally a single
equation of the form
\eqref{eq:twoCW} can balance
the ambient value
\(\lCW(Y)\) against a
nonzero derivative.
Two distinct ratios remove
that possibility. Neither
\(\lCW(Y)=0\) by itself nor
\(\Delta_K=1\) for an
arbitrary knot would
suffice for unknot detection.

\subsection{Generalized solid tori with torsion}

For odd \(g>1\), consider
the two-string tangle
\[
 T_g=\tang{1/g}+\tang{-1/g}.
\]
Its numerator closure is
the two-component unlink:
the two rational blocks
in that closure cancel.
Its denominator closure
is the connected sum of
the two opposite
\((2,g)\)-torus knots,
of determinant \(g^2\).
Its double branched cover
is the Seifert manifold
over a disk with two
exceptional fibres of
order \(g\), not a solid
torus.

One homology presentation
of this cover has
generators \(x,y,h\) and
relations
\(gx+h=0,\ gy-h=0\).
Eliminating \(h\) leaves
\(g(x+y)=0\), so its
homology is \(\Z\oplus
\Z/g\Z\). This is precisely the manifold \(N_g=M(\varnothing;1/g,-1/g)\)
of Rasmussen and Rasmussen. These are the
torsion-bearing examples
appearing among
generalized solid tori
in \cite[Section~7]{RR}.
They show concretely why
an unlink closure alone
does not force rationality.
They are excluded here
by the determinant-one
calibration giving
\(H_1(M)=\Z\).

\begin{remark}
This example also prevents
using an asserted
uniqueness of a
two-string solution to
one unlink-closure
equation as a shortcut.
In particular we do not
use Proposition~10(1)
of the arXiv version of
\cite{Sikora}. For
\(T_g\), the numerator
closure is the unlink
but the cover and the
denominator closure
distinguish it from the
trivial rational tangle.
The longitudinal case
in this paper instead
uses the homology
calibration and
\cite[Proposition~7.8]{RR}.
\end{remark}

\subsection{Orientation and marking}

An unoriented
homeomorphism of
three-manifolds does
not equate their
Casson--Walker
invariants: orientation
reversal changes its
sign. Thus a relation
\(M(\alpha)\cong
-V(\alpha)\) cannot
be substituted in
\eqref{eq:twoCW}.
Likewise knowing only
that two fillings are
lens spaces of the
same orders does not
identify their signed
lens-space terms.
The proof obtains
orientation-preserving
maps from actual
ambient link isotopies,
and obtains coefficients
from the common
boundary marking,
independently of those
maps.

\subsection{Local knotting and scalar equality}

Suppose hypothetically
that a nontrivial knot
\(J\) has \(V_J=1\).
Tie \(J\) locally into
one strand of a
rational tangle \(R\),
and correct the local
framing to make its
reduced bracket one.
The resulting tangle
\(X\) has \([X]=[R]\).
Every rational closure
has the same bracket
and determinant as
the corresponding
closure of \(R\).
Nevertheless the
closure that was an
unknot for \(R\) is
now \(J\), since that
closure traverses
the knotted strand.
An actual closure
identification is
therefore absent.
This hypothetical
example is a test of
the proof's logic,
not an assumption
or a construction
of a Jones-unit knot.

\section{What the result says about Jones unknot detection}\label{sec:jones}

Theorems~\ref{thm:framed}
and \ref{thm:three}
convert topological
closure information
into rationality.
The two actual
closures provide
the filling
homeomorphisms.
The scalar
information fixes
which peripheral
homology data to
use in those
homeomorphisms.
These roles cannot
be interchanged.

If a tangle has
the bracket vector
of a rational
tangle, all
rational closure
determinants
already agree.
The interpolation
lemma then
recovers the
integral homology
and longitude,
but it still
does not supply
a second actual
lens filling.
In particular,
an actual unknot
closure together
with the full
bracket vector
does not meet
the two-filling
hypothesis without
additional topology.

Sikora's implication
from Jones unknot
detection to framed
rational-solution
uniqueness therefore
remains a one-way
implication in this
discussion. Proving
its consequent
does not prove its
antecedent.
The stronger
assertion that the
bracket detects
rational tangles
among arbitrary
tangles is not
proved here.
The local-knotting
test in
Section~\ref{sec:stress}
pinpoints the
missing actual
closure hypotheses.

\section{Prior work, attribution, and verification}\label{sec:credit}

\subsection{The method antecedent}

Gainullin's
\cite[Theorem~8.2]{Gainullin}
characterizes
unknotted
null-homologous
knots in an
ambient \(L\)-space
under an appropriate
graded Heegaard
Floer comparison
with surgery on
the unknot.
Its proof combines
a surgery comparison,
the Casson--Walker
formula, a vanishing
second Alexander
derivative, the
alternating
\(L\)-space
Alexander form,
and genus detection.
That chain is the
antecedent for
Step~3 here.

There are four
specific differences
in the present
application.
Our reference
filling is
\(L(p,q)\), not
\(Y\#L(p,q)\).
The second matched
slope eliminates
\(\lCW(Y)\).
Floer simplicity
of the exterior
supplies the
torsion normal
form, without an
ambient \(L\)-space
hypothesis on \(Y\).
Finally
\cite[Corollary~2.3]{RR}
and irreducibility
turn degree zero
into a solid
torus via boundary
compression.
Choosing a lens
filling as the
ambient space
would not remove
these distinctions:
its core can be
homologically
essential, whereas
Gainullin's
theorem concerns
null-homologous
knots.

\subsection{Classical special cases and tangle rigidity}

Kronheimer,
Mrowka, Ozsv\'ath
and Szab\'o
\cite[Theorem~1.1]{KMOS}
characterize the
unknot in \(S^3\)
by a matched
oriented rational
surgery. If an
actual closure
in our problem
is an unknot,
its cover can be
used as \(S^3\)
after the
peripheral
calibration,
and this is a
classical special
case. We do not
present it as
new surgery
rigidity.

Ernst and
Sumners
\cite{ES}
developed
two-closure
rational-tangle
methods in the
DNA setting.
Buck and
Mauricio
\cite{BM}
and Baker and
Buck
\cite{BakerBuck}
study closely
related tangle
equations and
rational
subtangle
replacements.
These are
antecedents to
the rigidity
questions, not
claims that
two unframed
closures always
determine a
tangle. The
example in
Section~\ref{subsec:ES}
shows that they
do not.
The three-probe
theorem adds one
scalar measurement
and explains
exactly how it
removes the
homological
reflection
\eqref{eq:reflected}.

The bounded
literature review
for this draft
did not identify
an earlier proof
of the full
framed conjecture
or the exact
marked
three-probe
statement.
Sikora's October
2026 talk
abstract
\cite{SikoraTalk}
still describes
the former as
a conjecture.
The abstract,
not the complete
recording, was
checked. No
priority claim
is inferred
from that
observation.

\appendix
\section{Source versions and logical dependencies}

The following map records the external statements on which the deduction
rests. It is also intended to distinguish results about actual topology
from results about scalar invariants. All references to Sikora and to
Rasmussen--Rasmussen use the arXiv versions fixed in the bibliography.

\begin{center}
\small
\begin{tabular}{>{\raggedright\arraybackslash}p{0.23\textwidth}>{\raggedright\arraybackslash}p{0.29\textwidth}>{\raggedright\arraybackslash}p{0.36\textwidth}}
\hline
Input & Checked location & Use in this paper \\
\hline
Two probes determine the bracket vector & Sikora, Theorem 21 (journal 8.1) & Step 1; only framed closure hypotheses \\
Rational framed bracket injectivity & Sikora, Corollary 20 (journal 7.2) & Step 5; only after rationality is proved \\
Casson--Walker surgery & Ni--Wu, Proposition 3.1, (7) & Step 3; null-homologous core in the auxiliary integral homology sphere \\
Two \(L\)-space slopes & RR, Proposition 1.3 & Floer simplicity of the exterior, not of the auxiliary ambient manifold \\
Binary torsion & RR, paragraph before Proposition 1.4, Section 2.1, Propositions 1.4 and 2.2 & Alternation and positivity in Step 3 \\
Thurston norm & RR, Corollary 2.3, first assertion & Contradiction to boundary incompressibility in Steps 3 and 4 \\
Longitudinal filling & RR, Proposition 7.8 and Definition 7.2 & Step 4 only; \(Z=S^3\), \(g_M=1\) \\
Rationality of the given quotient & Buck--Mauricio, Theorem 3.5(2), after Lickorish & Step 5 and the unframed theorem \\
\hline
\end{tabular}
\end{center}

The unframed theorem replaces the first row by the elementary homology
interpolation argument. It then uses the nonlongitudinal surgery proof
and rational-tangle slope classification. It does not use RR Proposition
7.8 or any assertion that a bracket vector detects a rational tangle among
all tangles. The arithmetic criterion uses the same surgery proof after
eliminating the second integral functional in \eqref{eq:reflected}.

The title of the Buck--Mauricio article changed between preprint and
journal publication. Its result numbers here are pinned to the preprint.
The Boyer--Lines surgery formula and Lickorish's branched-cover criterion
are used through the explicitly identified formulations in Ni--Wu and
Buck--Mauricio; we do not assign unchecked original theorem numbers to
those two classical sources.

\section*{Statement on the use of AI systems}
The mathematics in this paper was developed in a collaboration directed by the author with two AI systems: a GPT-6 instance (OpenAI, through Codex) and Claude Fable 5.1 (Anthropic, through Claude Code). The proof of the framed theorem was produced by the GPT-6 instance on 1 November 2026 and audited independently by both systems in several rounds; the unframed extension was produced by the GPT-6 instance and checked by Claude on 3 November 2026. The author takes responsibility for every statement in the paper. The argument has not yet received external specialist review, and the priority of the result has not been independently established.

\begin{thebibliography}{99}

\bibitem{BakerBuck}
K.~L. Baker and D.~Buck,
\emph{The classification of rational subtangle replacements
between rational tangles},
Algebr. Geom. Topol. \textbf{13} (2013), 1413--1463.
\href{https://arxiv.org/abs/1304.7669}{arXiv:1304.7669}.
\href{https://doi.org/10.2140/agt.2013.13.1413}{doi:10.2140/agt.2013.13.1413}.

\bibitem{BoyerLines}
S.~Boyer and D.~Lines,
\emph{Surgery formulae for Casson's invariant and extensions
to homology lens spaces},
J. Reine Angew. Math. \textbf{405} (1990), 181--220.
The surgery formula used here is cited in the explicit
formulation of \cite[Proposition~3.1]{NiWu}.

\bibitem{BM}
D.~Buck and M.~Mauricio,
\emph{Connect sum of lens spaces surgeries: application
to Hin recombination},
Math. Proc. Cambridge Philos. Soc. \textbf{150} (2011),
505--525.
\href{https://arxiv.org/abs/1007.0948}{arXiv:1007.0948}, with preprint title
\emph{Tangle solutions for composite knots: applications to Hin recombination}.
Theorem numbers~3.5(2) and~3.6 refer to the inspected preprint.
\href{https://doi.org/10.1017/S0305004111000090}{doi:10.1017/S0305004111000090}.

\bibitem{DFKLS}
O.~T. Dasbach, D.~Futer, E.~Kalfagianni,
X.-S. Lin, and N.~W. Stoltzfus,
\emph{Alternating sum formulae for the determinant
and other link invariants},
J. Knot Theory Ramifications \textbf{19} (2010),
765--782.
\href{https://arxiv.org/abs/math/0611025}{arXiv:math/0611025}.

\bibitem{ES}
C.~Ernst and D.~W. Sumners,
\emph{A calculus for rational tangles: applications
to DNA recombination},
Math. Proc. Cambridge Philos. Soc. \textbf{108}
(1990), 489--515.

\bibitem{ES1999}
C.~Ernst and D.~W. Sumners,
\emph{Solving tangle equations arising in a DNA recombination model},
Math. Proc. Cambridge Philos. Soc. \textbf{126} (1999), 23--36.
\href{https://doi.org/10.1017/S0305004198002989}{doi:10.1017/S0305004198002989}.

\bibitem{Gainullin}
F.~Gainullin,
\emph{Heegaard Floer homology and knots determined
by their complements},
Algebr. Geom. Topol. \textbf{18} (2018), 69--109.
\href{https://arxiv.org/abs/1504.06180}{arXiv:1504.06180}.
\href{https://doi.org/10.2140/agt.2018.18.69}{doi:10.2140/agt.2018.18.69}.

\bibitem{KMOS}
P.~Kronheimer, T.~Mrowka, P.~Ozsv\'ath, and Z.~Szab\'o,
\emph{Monopoles and lens space surgeries},
Ann. of Math. (2) \textbf{165} (2007), 457--546.
\href{https://arxiv.org/abs/math/0310164}{arXiv:math/0310164}.
\href{https://doi.org/10.4007/annals.2007.165.457}{doi:10.4007/annals.2007.165.457}.

\bibitem{Lickorish}
W.~B.~R. Lickorish,
\emph{Prime knots and tangles},
Trans. Amer. Math. Soc. \textbf{267} (1981),
321--332.
The branched-cover criterion is used through
\cite[Theorem~3.5(2)]{BM}.

\bibitem{LickorishBook}
W.~B.~R. Lickorish,
\emph{An Introduction to Knot Theory},
Graduate Texts in Mathematics 175, Springer, New York, 1997.
Chapter~9, Theorems~9.3--9.4 and Corollary~9.5.
\href{https://doi.org/10.1007/978-1-4612-0691-0_9}{doi:10.1007/978-1-4612-0691-0\_9}.

\bibitem{NiWu}
Y.~Ni and Z.~Wu,
\emph{Cosmetic surgeries on knots in \(S^3\)},
J. Reine Angew. Math. \textbf{706} (2015), 1--17.
\href{https://arxiv.org/abs/1009.4720}{arXiv:1009.4720v2}.
\href{https://doi.org/10.1515/crelle-2013-0067}{doi:10.1515/crelle-2013-0067}.

\bibitem{RR}
J.~Rasmussen and S.~D. Rasmussen,
\emph{Floer simple manifolds and \(L\)-space intervals},
Adv. Math. \textbf{322} (2017), 738--805.
\href{https://arxiv.org/abs/1508.05900}{arXiv:1508.05900v2}.
\href{https://doi.org/10.1016/j.aim.2017.10.014}{doi:10.1016/j.aim.2017.10.014}.
Numbered results in this manuscript refer to this
specified arXiv version.

\bibitem{Sikora}
A.~S. Sikora,
\emph{Tangle equations, the Jones conjecture, slopes of
surfaces in tangle complements, and \(q\)-deformed rationals},
Canad. J. Math. \textbf{76} (2024), 707--727.
\href{https://arxiv.org/abs/2005.08162}{arXiv:2005.08162v2}.
\href{https://doi.org/10.4153/S0008414X23000755}{doi:10.4153/S0008414X23000755}.
Unless explicitly stated otherwise, numbered results
refer to arXiv v2: Proposition~1, Corollary~20, and
Theorem~21 are journal Proposition~1.1, Corollary~7.2,
and Theorem~8.1, respectively.

\bibitem{SikoraTalk}
A.~S. Sikora,
\emph{Tangle equations, the Jones unknot conjecture,
and lens space surgeries},
IPAM, Workshop I: Hyperbolic Structures and Quantum
Invariants, 15 October 2026.
Talk abstract,
\url{https://www.mathinstitutes.org/videos/25091}
(accessed 3 November 2026).

\bibitem{Stoimenow}
A.~Stoimenow,
\emph{Square numbers, spanning trees and invariants
of achiral knots},
Comm. Anal. Geom. \textbf{13} (2005), 591--631.
\href{https://doi.org/10.4310/CAG.2005.v13.n3.a5}{doi:10.4310/CAG.2005.v13.n3.a5}.

\end{thebibliography}
\end{document}
